(1) 对不透明漫灰板，离开表面的辐射度满足 $J_i=\epsilon_i\sigma_{\rm SB}T_i^4+(1-\epsilon_i)G_i$。无限平板视角因子为一，所以 $G_1=J_2,G_2=J_1$，两板之间净热流为 $q=J_1-J_2$，正号取从高温板到低温板。
(2) 将 $G_1=J_2,G_2=J_1$ 分别代回辐射度式，得 $J_1-\sigma_{\rm SB}T_1^4=-(1-\epsilon_1)q/\epsilon_1$、$J_2-\sigma_{\rm SB}T_2^4=(1-\epsilon_2)q/\epsilon_2$。消去 $J_i$ 得 $q=\sigma_{\rm SB}(T_1^4-T_2^4)/(1/\epsilon_1+1/\epsilon_2-1)$。
(3) 若 $\epsilon_1=\epsilon_2=1$，分母为一，得黑体换热 $q=\sigma_{\rm SB}(T_1^4-T_2^4)$。若任何一侧 $\epsilon_i\to0^+$，分母趋无穷，净交换趋零；这体现理想全反射表面无法直接有效辐射换热。
(4) 设两板温度接近共同温度 $T$，使用 $T_1^4-T_2^4\simeq4T^3(T_1-T_2)$，因此 $q\simeq h_r(T_1-T_2)$，$h_r=4\sigma_{\rm SB}T^3/(1/\epsilon_1+1/\epsilon_2-1)$。其单位为 $\mathrm{W\,m^{-2}\,K^{-1}}$，符合表面换热系数。
